% ==========================================
% sec1_1.tex 开头：局部定义随机变量（一维、二维）相关引理颜色
% ==========================================
\makeatletter
% 定义 sec1_1 的专属颜色：深蓝色
\definecolor{LCCColorSec11}{RGB}{20, 40, 100} 
% 定义新的引理环境，继承 lcc@basestyle 和原有的计数器
\newtcolorbox[use counter from=引理]{引理Sec11}[1][]{
    lcc@basestyle,
    colframe=LCCColorSec11!70!white, colback=LCCColorSec11!2,
    interior style={top color=LCCColorSec11!8, bottom color=LCCColorSec11!1},
    frame style={left color=LCCColorSec11!70!white, right color=LCCDeepBlue},
    title={引理~\thetcbcounter\if\relax\detokenize{#1}\relax\else\quad #1\fi}
}
\makeatother
% ------------------------------------------



% ==========================================
% 离散型随机变量及其分布
% ==========================================

\subsubsection*{1. 0-1 分布与二项分布}

\textbf{0-1 分布（特殊的二项分布）}
\begin{itemize}
    \item 期望：$E(X) = p$
    \item 方差：$D(X) = p(1-p)$
\end{itemize}

\vspace{0.5em}
\textbf{二项分布} \quad $X \sim b(n,p)$ 或 $X \sim B(n,p)$
\begin{itemize}
    \item 分布律：$P(X=k) = C_n^k p^k (1-p)^{n-k}, \quad k=0,1,2,\cdots,n$
    \item 期望：$E(X) = np$
    \item 方差：$D(X) = np(1-p)$
\end{itemize}
\textbf{证明提示}\quad 
将 $X$ 看成 $n$ 个独立同分布的 0-1 分布变量 $X_i$ 的和，即 $X = \sum_{i=1}^n X_i$。由于期望和方差（独立前提下）具有可加性，因此 $E(X) = \sum E(X_i) = np$，且 $D(X) = \sum D(X_i) = np(1-p)$。

\vspace{1em}
\subsubsection*{2. 泊松分布}

\begin{引理Sec11}[指数函数的麦克劳林展开]
	对于任意实数 $\lambda$，指数函数 $e^\lambda$ 可以展开为幂级数，这是推导泊松分布数字特征的核心基础：
	$$ e^\lambda = \sum_{k=0}^{\infty} \frac{\lambda^k}{k!} = 1 + \lambda + \frac{\lambda^2}{2!} + \cdots $$
\end{引理Sec11}

$X \sim \pi(\lambda)$ 或 $X \sim P(\lambda)$
\begin{itemize}
    \item 分布律：$P(X=k) = \frac{\lambda^k}{k!} e^{-\lambda}, \quad k=0,1,2,\cdots$
    \item 泊松定理（二项分布的极限）：当 $n \ge 20, p \le 0.05$ 时，有近似公式 $C_n^k p^k (1-p)^{n-k} \approx \frac{\lambda^k}{k!} e^{-\lambda}$，其中 $\lambda = np$。
\end{itemize}

\textbf{期望与方差的证明}\quad

首先求期望 $E(X)$：
$$ E(X) = \sum_{k=0}^{\infty} k \cdot \frac{\lambda^k}{k!} e^{-\lambda} = e^{-\lambda} \sum_{k=1}^{\infty} k \cdot \frac{\lambda^k}{k(k-1)!} = e^{-\lambda} \cdot \lambda \sum_{k=1}^{\infty} \frac{\lambda^{k-1}}{(k-1)!} = e^{-\lambda} \cdot \lambda \cdot e^{\lambda} = \lambda $$

接着求 $E(X^2)$，这里的核心技巧是利用 $k^2 = k(k-1) + k$ 进行拆项：
$$ 
\begin{aligned}
E(X^2) &= \sum_{k=0}^{\infty} k^2 \frac{\lambda^k}{k!} e^{-\lambda} = \sum_{k=1}^{\infty} (k(k-1) + k) \frac{\lambda^k}{k!} e^{-\lambda} \\
&= \sum_{k=2}^{\infty} k(k-1) \frac{\lambda^k}{k!} e^{-\lambda} + \sum_{k=1}^{\infty} k \frac{\lambda^k}{k!} e^{-\lambda} \\
&= \lambda^2 e^{-\lambda} \sum_{k=2}^{\infty} \frac{\lambda^{k-2}}{(k-2)!} + E(X) \\
&= \lambda^2 e^{-\lambda} \cdot e^{\lambda} + \lambda = \lambda^2 + \lambda
\end{aligned}
$$
最后计算方差：
$$ D(X) = E(X^2) - [E(X)]^2 = (\lambda^2 + \lambda) - \lambda^2 = \lambda $$


\vspace{1em}
\subsubsection*{3. 几何分布与超几何分布}

\begin{引理Sec11}[幂级数逐项求导]
	对于无穷等比级数 $\sum_{k=1}^{\infty} q^k = \frac{q}{1-q}$ （其中 $|q| < 1$），在收敛域内可以逐项求导。即：
	$$ \sum_{k=1}^{\infty} k q^{k-1} = \frac{1}{(1-q)^2} $$
	同理，对上式再次求导可得：
	$$ \sum_{k=2}^{\infty} k(k-1) q^{k-2} = \frac{2}{(1-q)^3} $$
\end{引理Sec11}

\textbf{几何分布}
\begin{itemize}
    \item 分布律：$P(X=k) = (1-p)^{k-1}p, \quad k=1,2,3,\cdots$
    \item 期望与方差：$E(X) = \frac{1}{p}, \quad D(X) = \frac{1-p}{p^2}$
\end{itemize}

\textbf{期望与方差的证明}\quad 

令 $q = 1-p$。利用上述引理，直接代入可得期望：
$$ E(X) = \sum_{k=1}^{\infty} k \cdot q^{k-1}p = p \sum_{k=1}^{\infty} k q^{k-1} = p \cdot \frac{1}{(1-q)^2} = p \cdot \frac{1}{p^2} = \frac{1}{p} $$
求方差时，先求 $E[X(X-1)]$：
$$ E[X(X-1)] = \sum_{k=1}^{\infty} k(k-1) q^{k-1}p = pq \sum_{k=2}^{\infty} k(k-1)q^{k-2} = pq \cdot \frac{2}{(1-q)^3} = pq \cdot \frac{2}{p^3} = \frac{2q}{p^2} $$
$$ D(X) = E[X(X-1)] + E(X) - [E(X)]^2 = \frac{2(1-p)}{p^2} + \frac{1}{p} - \frac{1}{p^2} = \frac{1-p}{p^2} $$

\vspace{1em}
\begin{引理Sec11}[组合数恒等式]
	在组合数学中，有以下两个常用于概率论推导的重要恒等式：
	\begin{itemize}
		\item 吸收恒等式：$k C_M^k = M C_{M-1}^{k-1}$，同理 $k(k-1)C_M^k = M(M-1)C_{M-2}^{k-2}$。
		\item 范德蒙德卷积公式：$\sum_{k} C_A^k C_B^{n-k} = C_{A+B}^n$。
	\end{itemize}
\end{引理Sec11}

\textbf{超几何分布} \quad 场景：$N$ 件物品，其中 $M$ 件次品，不放回地抽取 $n$ 件。$X$ 表示抽到的次品数。
\begin{itemize}
    \item 分布律：$P(X=k) = \frac{C_M^k \cdot C_{N-M}^{n-k}}{C_N^n}, \quad k=0,1,2,\cdots,\min\{n, M\}$
    \item 期望：$E(X) = np$，其中 $p = \frac{M}{N}$ 为次品率。
    \item 方差：$D(X) = np(1-p)\frac{N-n}{N-1}$，其中 $p = \frac{M}{N}$ 为次品率。
\end{itemize}

\textbf{期望与方差的证明}\quad 

证明的核心在于利用引理中的吸收恒等式以及范德蒙德卷积公式。
$$ 
\begin{aligned}
E(X) &= \sum_{k=1}^{n} k \frac{C_M^k C_{N-M}^{n-k}}{C_N^n} = \sum_{k=1}^{n} \frac{M C_{M-1}^{k-1} C_{N-M}^{n-k}}{\frac{N}{n} C_{N-1}^{n-1}} \\
&= n \frac{M}{N} \sum_{k=1}^{n} \frac{C_{M-1}^{k-1} C_{(N-1)-(M-1)}^{(n-1)-(k-1)}}{C_{N-1}^{n-1}}
\end{aligned}
$$
由于后半部分的求和代表了在 $N-1$ 件物品中抽取 $n-1$ 件的所有概率之和，其值等于 $1$。因此：
$$ E(X) = n \frac{M}{N} = np $$

求方差同样先求 $E[X(X-1)]$：
$$ 
\begin{aligned}
E[X(X-1)] &= \sum_{k=2}^{n} k(k-1) \frac{C_M^k C_{N-M}^{n-k}}{C_N^n} = \sum_{k=2}^{n} \frac{M(M-1) C_{M-2}^{k-2} C_{N-M}^{n-k}}{\frac{N(N-1)}{n(n-1)} C_{N-2}^{n-2}} \\
&= \frac{n(n-1)M(M-1)}{N(N-1)} \sum_{k=2}^{n} \frac{C_{M-2}^{k-2} C_{(N-2)-(M-2)}^{(n-2)-(k-2)}}{C_{N-2}^{n-2}} \\
&= \frac{n(n-1)M(M-1)}{N(N-1)}
\end{aligned}
$$
最后代入方差公式化简（令 $p = \frac{M}{N}$）：
$$ 
\begin{aligned}
D(X) &= E[X(X-1)] + E(X) - [E(X)]^2 \\
&= \frac{n(n-1)M(M-1)}{N(N-1)} + np - n^2 p^2 \\
&= np \left[ \frac{(n-1)(M-1)}{N-1} + 1 - np \right] \\
&= np(1-p)\frac{N-n}{N-1}
\end{aligned}
$$


% ==========================================
% 连续型随机变量及其分布
% ==========================================
\vspace{1.5em}

\subsubsection*{4. 指数分布}
\begin{引理Sec11}[分部积分法]
	在计算连续型随机变量（如指数分布、正态分布）的数字特征时，常利用分部积分法消去被积函数中的多项式因子：
	$$ \int u \, dv = uv - \int v \, du $$
\end{引理Sec11}
$X \sim E(\lambda)$

也有参数化形式为 $\theta = \frac{1}{\lambda}$，请根据任课老师要求记忆：
\begin{itemize}
    \item 概率密度函数：$f(x) = \begin{cases} \frac{1}{\theta} e^{-\frac{x}{\theta}}, & x > 0 \\ 0, & \text{其他} \end{cases}$
    \item 分布函数：$F(x) = \begin{cases} 1 - e^{-\frac{x}{\theta}}, & x > 0 \\ 0, & \text{其他} \end{cases}$
\end{itemize}

\textbf{无记忆性证明}\quad ，即证明 $P\{X > s+t \mid X > s\} = P\{X > t\}$：
$$ 
\begin{aligned}
P\{X > s+t \mid X > s\} &= \frac{P\{X > s+t \text{ 且 } X > s\}}{P\{X > s\}} = \frac{P\{X > s+t\}}{P\{X > s\}} \\
&= \frac{1 - F(s+t)}{1 - F(s)} = \frac{e^{-\frac{s+t}{\theta}}}{e^{-\frac{s}{\theta}}} = e^{-\frac{t}{\theta}} \\
&= 1 - F(t) = P\{X > t\}
\end{aligned}
$$

\textbf{期望与方差的证明（利用分部积分法）}\quad
$$ 
\begin{aligned}
E(X) &= \int_{0}^{+\infty} x \cdot \frac{1}{\theta} e^{-\frac{x}{\theta}} dx = -\int_{0}^{+\infty} x \, d(e^{-\frac{x}{\theta}}) \\
&= -\left[ x e^{-\frac{x}{\theta}} \right]_{0}^{+\infty} + \int_{0}^{+\infty} e^{-\frac{x}{\theta}} dx \\
&= 0 + \left[ -\theta e^{-\frac{x}{\theta}} \right]_{0}^{+\infty} = \theta
\end{aligned}
$$
同理，两次运用分部积分法求 $E(X^2)$：
$$ E(X^2) = \int_{0}^{+\infty} x^2 \cdot \frac{1}{\theta} e^{-\frac{x}{\theta}} dx = 2\theta^2 $$
$$ D(X) = E(X^2) - [E(X)]^2 = 2\theta^2 - \theta^2 = \theta^2 $$


\vspace{1em}
\subsubsection*{5. 均匀分布}
$X \sim U(a,b)$
\begin{itemize}
    \item 概率密度函数：$f(x) = \begin{cases} \frac{1}{b-a}, & a < x < b \\ 0, & \text{其他} \end{cases}$
\end{itemize}

\textbf{期望与方差的证明}\quad 
$$ E(X) = \int_{a}^{b} x \cdot \frac{1}{b-a} dx = \frac{1}{b-a} \left[ \frac{1}{2}x^2 \right]_{a}^{b} = \frac{b^2 - a^2}{2(b-a)} = \frac{a+b}{2} $$
求方差先求 $E(X^2)$：
$$ E(X^2) = \int_{a}^{b} x^2 \cdot \frac{1}{b-a} dx = \frac{1}{b-a} \left[ \frac{1}{3}x^3 \right]_{a}^{b} = \frac{b^3 - a^3}{3(b-a)} = \frac{1}{3}(a^2 + ab + b^2) $$
$$ 
\begin{aligned}
D(X) &= E(X^2) - [E(X)]^2 = \frac{a^2 + ab + b^2}{3} - \frac{(a+b)^2}{4} \\
&= \frac{4a^2 + 4ab + 4b^2 - 3(a^2 + 2ab + b^2)}{12} \\
&= \frac{a^2 - 2ab + b^2}{12} = \frac{1}{12}(b-a)^2
\end{aligned}
$$


\vspace{1em}
\subsubsection*{6. 正态分布}
$X \sim N(\mu, \sigma^2)$
\begin{itemize}
    \item 概率密度函数：$f(x) = \frac{1}{\sqrt{2\pi}\sigma} e^{-\frac{(x-\mu)^2}{2\sigma^2}}$
    \item 期望与方差：$E(X) = \mu, \quad D(X) = \sigma^2$
\end{itemize}

\textbf{重要性质}
\begin{itemize}
    \item 标准正态分布对称性：若 $X \sim N(0,1)$，则 $\Phi(x) + \Phi(-x) = 1$。
    \item 标准化变换：若 $X \sim N(\mu, \sigma^2)$，则 $\frac{X-\mu}{\sigma} \sim N(0,1)$。
\end{itemize}